\section{Verification purpose}

This case verifies the effect of nonzero surface spread on ignition synchronization in a fine-grid WUI calculation. A no-surface reference and a surface-coupled simulation condition use identical geometry, wind, firebrand physics, and resolution; only surface propagation changes.

The expected effect is earlier and more uniform ignition within each structure while cell-scale firebrand-accumulation distortion remains measurable. Success is based on paired structure-level timing and load metrics rather than visual comparison alone.

\section{Mathematical formulation and numerical implementation}

\subsection{Fine-grid structural discretization}

At \(\Delta x=\SI{2.5}{m}\), each \SI{10}{m} structure is represented by
a \(4\times4\) block of square cells. For cell \(j\), the
structure heat-release rate (HRR) used by the firebrand model is
\[
  \mathrm{HRR}_j(t)=q''(t)\,\Delta x^2.
\]
No strip-width correction is applied: summing the 16 cells recovers the
physical \SI{100}{m^2} footprint and its design HRR. Firebrands are generated with
\[
  \dot N_j(t)=GR'\,\mathrm{HRR}_j(t),\qquad
  GR'=\SI{10}{pcs.MW^{-1}.s^{-1}},
\]
and transported by the current Eulerian implementation using the structural
Himoto distribution. Physical ignition uses the critical accumulated load
\(N''_{\rm crit}=\SI{2928.9}{pcs.m^{-2}}\), followed by the configured
small-fuel and development delays.

In the no-surface simulation condition, cell-local deposition makes downwind cells within
one structure
receive different loads. Those cells consequently cross the ignition
criterion at different times, producing artificial intra-structure delay and
slow community-scale propagation.

\subsection{Surface-fire synchronization mechanism}

The surface-coupled simulation condition retains the firebrand process but activates surface
spread throughout the same two-dimensional lattice. The configured Hamada
surface rate is
\(ROS_{\rm surface}\approx\SI{0.34}{m.s^{-1}}\)~\cite{Qin2025}. Once one cell of a
\SI{10}{m} structure ignites, a surface front at this speed traverses the
footprint in approximately
\[
  \Delta t_{\rm traverse}
  \approx \frac{a}{ROS_{\rm surface}}
  =\frac{\SI{10}{m}}{\SI{0.34}{m.s^{-1}}}
  \approx\SI{29.4}{s}.
\]
This is much shorter than the cell-to-cell ignition offsets created by
firebrand accumulation alone. Surface coupling should therefore reduce
within-structure ignition spans and accelerate the downwind ignition
chronology. It does not alter the cell-local firebrand deposition algorithm,
so load imbalance should remain visible.

\section{Derivation of expected behavior}
Surface propagation crosses a represented structure much faster than independent firebrand accumulation ignites its cells. The coupled simulation condition should therefore ignite earlier and with a smaller within-structure time span, while the unchanged deposition algorithm should preserve measurable load nonuniformity.

% BEGIN EXPLICIT EXPECTED-RESULT DERIVATION
At \(\Delta x=2.5~\mathrm{m}\), each structural cell has area
\(2.5^2=6.25~\mathrm{m^2}\).  At the design-fire plateau,
\[
 HRR_{j,\max}=400(6.25)=2500~\mathrm{kW}=2.5~\mathrm{MW},
 \qquad \dot N_{j,\max}=10(2.5)=25~\mathrm{pcs\,s^{-1}}.
\]
Summing the \(4\times4=16\) cells gives
\(16(2.5)=40~\mathrm{MW}\) and
\(16(25)=400~\mathrm{pcs\,s^{-1}}\), exactly the physical
\(100~\mathrm{m^2}\) structure totals.  The generated level-set step is
\[
 \Delta t=0.5(2.5)/[40(0.44704)]=0.0699043~\mathrm{s}.
\]
Once one structural cell ignites, the surface reference crosses one 2.5-m
cell in \(2.5/0.34=7.3529~\mathrm{s}\) and a full 10-m footprint in
\(10/0.34=29.4118~\mathrm{s}\).  These are the explicit synchronization
timescales; firebrand-only cells still wait for independent accumulation to
\(2928.9~\mathrm{pcs\,m^{-2}}\) plus the configured ignition delays.

The ideal paired signatures are therefore
\(t_{\mathrm{coupled}}<t_{\mathrm{no\ surface}}\) for most matched cells,
a smaller coupled within-structure span, and a larger coupled fitted rate of spread (ROS),
while the deposition range remains nonzero because its algorithm is unchanged.
The quantified substitutions require at least 75 percent earlier cells,
median time ratio at most 0.90, median span ratio at most 0.50, ROS ratio at
least 2, and coupled load range at least 0.05.
% END EXPLICIT EXPECTED-RESULT DERIVATION

\section{Verification metrics and acceptance criteria}

For every downstream structural cell having finite arrival times in both
simulation conditions, let \(t_{\mathrm{ns},j}\) denote the no-surface arrival and
\(t_{\mathrm{sc},j}\) the surface-coupled arrival. Earlier ignition is
quantified by
\[
  F_{\rm early}=\frac{1}{N}\sum_j
  \mathbf{1}(t_{\mathrm{sc},j}<t_{\mathrm{ns},j}),
  \qquad
  R_t=\operatorname{median}_j\left(
  \frac{t_{\mathrm{sc},j}}{t_{\mathrm{ns},j}}\right).
\]
The case requires \(F_{\rm early}\ge0.75\) and \(R_t\le0.90\).

For structure \(s\), define its ignition span as
\[
  S_s=\max_{j\in s}t_j-\min_{j\in s}t_j.
\]
The synchronization metric is the median paired ratio
\[
  R_S=\operatorname{median}_{s\ge2}
      \left(\frac{S_{\mathrm{sc},s}}{S_{\mathrm{ns},s}}\right),
\]
which must not exceed 0.50. Mean propagation speed is obtained from a
least-squares fit of first-ignition time against structure center,
\(ROS=(dt/dx)^{-1}\); the coupled/no-surface ratio must be at least 2.

Persistence of accumulation distortion is measured from the represented cell
loads along the diagnostic lattice row in each coupled structure:
\[
  D_s=\frac{\max_{j\in s}N''_j-\min_{j\in s}N''_j}
            {\operatorname{mean}_{j\in s}N''_j}.
\]
The downstream median must remain at or above 0.05. This is intentionally a
lower-bound test: the surface-coupled simulation condition is expected to improve ignition
synchronization,
not to homogenize firebrand deposition.

\subsection{Justification of metric quantification}

These are minimum effect-size criteria for the directional comparison reported
for this configuration~\cite{Qin2025}, not estimates of experimental
uncertainty.  Requiring \(F_{\rm early}\ge0.75\) means earlier ignition must
occur in a clear supermajority of matched cells, so a few ties or grid-crossing
reversals cannot determine the result.  The median time ratio of 0.90 requires
at least a 10\% central timing improvement, while the span ratio of 0.50
requires a factor-of-two synchronization improvement within structures.  The
ROS ratio of 2.0 likewise demands a large propagation effect consistent with
surface coupling, rather than declaring success from an arbitrarily small
positive change.

The 5\% lower bound on within-structure load range is deliberately above
floating-point and raster-write noise.  It confirms that the known deposition
nonuniformity remains visible even when ignition becomes synchronized; it is
not an accuracy tolerance.  Requiring both input-matched simulation conditions is
essential because all five quantities are paired effects and have no valid
unpaired interpretation.

\subsection{Predeclared acceptance criteria}

Table~\ref{tab:acceptance-contract} summarizes the metrics fixed before
inspection of the calculated results. Numerical values and component statuses
are reported separately in Section~7.

\begin{table}[H]
\centering
\footnotesize
\caption{Predeclared verification metrics and acceptance criteria.}
\label{tab:acceptance-contract}
\begin{tabular}{@{}p{0.23\linewidth}p{0.47\linewidth}p{0.21\linewidth}@{}}
\toprule
Metric & Output, region, and aggregation & Acceptance criterion\\
\midrule
Output completeness & Both simulation conditions must carry the current input identity record and usable structure-level load/ignition fields. & 2 of 2 simulation conditions \\
Earlier-ignition fraction & Fraction of matched downstream cells with \(t_{coupled}<t_{no\ surface}\). & \(\ge0.75\) \\
Median ignition-time ratio & Median \(t_{coupled}/t_{no\ surface}\) over matched downstream cells. & \(\le0.90\) \\
Ignition-span ratio & Median paired within-structure span ratio over downstream structures. & \(\le0.50\) \\
Mean-ROS ratio & Coupled fitted structure-propagation ROS divided by no-surface ROS. & \(\ge2.0\) \\
Persistent load distortion & Median coupled within-structure relative load range. & \(\ge0.05\) \\
Overall decision & Conjunction of completeness and all five comparison metrics. & PASS only if all pass \\
\bottomrule
\end{tabular}
\end{table}

\section{Simulation configuration and predicted outcome}

Figure~\ref{fig:input-configuration} records the prepared spatial inputs for the 2.5 m grid with surface spread enabled.
These panels show the prepared spatial fields, remove the
two-cell numerical buffer, and retain map coordinates in metres.  The left panel shows
the most informative fuel or structure layout available for the case; the right panel
shows the cells selected by the non-positive initial level-set field, making a localized
ignition or firebrand-emission source explicit.

\begin{figure}[H]
\centering
\EvidenceFigure[width=0.98\textwidth,height=0.42\textheight,keepaspectratio]{../figures/input_configuration.pdf}
\caption{Whole-domain prepared input configuration for the representative simulation condition.  The
physical domain is shown without the numerical halo; coordinates, configured wind values,
fuel or structure layout, and initial ignition/source geometry are identified in the panels.}
\label{fig:input-configuration}
\end{figure}

\begin{table}[H]
\centering\small
\caption{Controlled paired configuration.}
\begin{tabular}{lll}
\toprule
Parameter & No-surface simulation condition & Surface-coupled simulation condition\\
\midrule
Grid spacing & \SI{2.5}{m} & \SI{2.5}{m}\\
Physical transect & \SI{100}{m} by \SI{10}{m} & same\\
Structure/firebreak & \SI{10}{m}/\SI{10}{m} & same\\
Wind & \SI{40}{mph}, west to east & same\\
Time step & \(0.5\Delta x/U\) & same\\
Requested maximum time & \SI{6000}{s} & \SI{3500}{s}\\
Surface fuel & model 91 structures; model 102 gaps & same\\
Spread model & disabled & Hamada (a=d=\SI{10}{m})\\
Generation & per-MW & per-MW\\
Transport/accumulation & empirical/Eulerian & same\\
Ignition & physical & same\\
\bottomrule
\end{tabular}
\end{table}

The fixed wind Courant--Friedrichs--Lewy (CFL) requires a fractional step. Preprocessing sets each generated
stop to \(\lceil t_{\mathrm{requested}}/\Delta t\rceil\Delta t\), and records
the requested stop, aligned stop, and integer step count in the simulation specification.

The no-surface reference retains a longer observation window because
accumulation-only spread requires it to ignite multiple downstream structures.
The surface-coupled simulation condition stops
after all structures have ignited and a mature, nonuniform ember-load field
has formed. Since time of arrival (TOA) is immutable after ignition and load distortion is
measured within the coupled run rather than between terminal times, this
difference reduces execution without changing any acceptance quantity.

The expected signature is not agreement with a prescribed straight-line TOA
curve. It is a directional paired result: the surface-coupled simulation condition must
ignite earlier, have
smaller intra-structure ignition spans, and have a larger fitted mean ROS,
while its cell-level load curve remains nonuniform.

\subsection{Preprocessing, execution, and postprocessing}

The two prescribed conditions use spatially aligned fields in a Universal
Transverse Mercator coordinate system. Square structures and surface-fuel gaps
repeat in both horizontal directions over the \(100\times100\)-m physical
domain; only the two-cell halo uses inert fuel model 93. Structure properties
share that grid. ELMFIRE applies the Hamada velocity in structures and the
Rothermel velocity in surface-fuel gaps, providing a continuous two-dimensional
route for the coupled front. Every structure in the first upwind column is
initialized, and distinct structure identities preserve each physical
\(4\times4\)-cell footprint. The two conditions differ only in whether
vegetative surface spread is enabled. Input-identity checks exclude results
from earlier configurations.

The two conditions are simulated independently. Analysis uses their final
fire arrival times and accumulated firebrand counts, normalized by each cell's
physical \(\Delta x\,w_y\) area, to calculate the timing and load comparisons.

\section{Actual simulation results}

Figure~\ref{fig:domain-result} shows the representative time-of-arrival field from the 2.5 m grid with surface spread enabled.
The displayed field is taken from the simulated results, masks missing values, and excludes
the two-cell numerical buffer.  The map shows the spatial synchronization of surface spread and structure ignition that is reduced to profiles in the subsequent comparison.

\begin{figure}[H]
\centering
\EvidenceFigure[width=0.96\textwidth,height=0.50\textheight,keepaspectratio]{../figures/domain_result.pdf}
\caption{Whole-domain time-of-arrival field from the representative completed ELMFIRE run.  This
domain-scale result supplements, rather than replaces, the higher-level verification
profiles, convergence plots, ensemble statistics, and calculated acceptance metrics below.}
\label{fig:domain-result}
\end{figure}

\CompletedVariantCount{} of \RequestedVariantCount{} input-matched
simulation conditions have evaluable outputs. Figure~\ref{fig:wsc-reference} first presents
the surface-coupled experiment by itself using the two physical profiles that
define the expected synchronization behavior.

\begin{figure}[H]
\centering
\EvidenceFigure[width=\linewidth]{surface_coupled_reference.pdf}
\caption{Surface-coupled fine-grid result. Left: cell-level accumulated
firebrand load and the critical ignition load. Right: structural-cell ignition
chronology.}
\label{fig:wsc-reference}
\end{figure}

Figure~\ref{fig:wsc-comparison} then overlays the no-surface control, making the
earlier and more uniform coupled ignition response directly visible. The
overall decision is \textbf{\OverallStatus}.

\begin{figure}[H]
\centering
\EvidenceFigure[width=\linewidth]{surface_coupling_comparison.pdf}
\caption{Controlled surface-coupling comparison. Persistent within-structure
load variation diagnoses the remaining accumulation distortion, while the
lower, smoother coupled ignition curve demonstrates synchronization.}
\label{fig:wsc-comparison}
\end{figure}

\section{Calculated metrics and verification decision}

\begin{table}[H]
\centering\small
\caption{Calculated verification metrics and component decisions.}
\begin{tabular}{@{}p{0.37\linewidth}p{0.18\linewidth}p{0.17\linewidth}p{0.14\linewidth}@{}}
\toprule
Metric & Acceptance limit & Calculated & Status\\
\midrule
\MetricRows
\bottomrule
\end{tabular}
\end{table}

The current decision is \textbf{\OverallStatus}. A pass requires both
input-matched outputs and every synchronization, acceleration, and
persistent-distortion check in the table. Missing comparison values are
reported as \emph{NOT EVALUABLE}, never as a pass.

This is an implementation verification using an idealized two-dimensional WUI
lattice. The fitted ROS is a community-scale diagnostic and is not required
to equal the configured Hamada surface speed exactly because the executable
also retains spotting, Rothermel fuel behavior, and discrete ignition delays.
The case verifies the qualitative and quantitative coupling signature derived
above; it does not claim that surface spread corrects the underlying
cell-local firebrand transport distortion.

\section{Technical interpretation}
The no-surface reference completes at its aligned \SI{6000.0224}{s} stop and supplies a input-matched final output.  The surface-coupled member uses the predeclared single controlled difference, enabled vegetative surface spread, but does not complete its \SI{3500.0364}{s} run.

The coupled execution writes a valid first transient dump near \SI{500.025}{s}, then exhibits sustained memory growth and severe slowdown: resident memory rises from approximately 319 MB at 93 s wall time to 874 MB at 550 s, 1.07 GB at 823 s, and 1.28 GB at 1188 s without reaching the second \SI{1000}{s} dump.  It was stopped before exhausting the shared environment.  Because the preprocessor, input identity record, input rasters, and controlled namelist contrast are valid, changing the case to avoid the coupled path would invalidate the verification objective.  The report therefore records 1/2 completeness and \emph{NOT RUN}; the unavailable comparison metrics are not treated as passes.  The resource-growth behavior is an implementation issue in the coupled WUI execution path.

\section{Scope and limitations}
The paired comparison isolates one uniform fine-grid WUI transect. It verifies the coupling signature and does not establish that surface spread corrects the underlying deposition discretization or predicts a real community fire.

\begin{thebibliography}{9}
\bibitem{Qin2025}
Y. Qin, \emph{A Physics-Based Eulerian Framework for Modeling Firebrand
Showering in Regional-Scale Wildland and Wildland--Urban Interface Fire
Simulations}, Ph.D. dissertation, University of Maryland, College Park, 2025.
\end{thebibliography}
